\input{../tex-inputs/latex-leseansicht-vorspann}

\section[Arthur Schnitzler an Gustav Schwarzkopf, 20. 5. 1906]{L04037 Arthur Schnitzler an Gustav Schwarzkopf, 20. 5. 1906}
\nopagebreak\mylabel{L04037v}
\rehead{ }\normalsize\beginnumbering\briefempfaengerindex{Schwarzkopf, Gustav@\textsc{Schwarzkopf, Gustav}!zzzSchnitzler, Arthur@\emph{von Arthur Schnitzler}!1906-05-202@{20. 5. 1906}|(be}
\toendnotes[C]{\smallbreak\pagebreak[2]}
\correspDesc{Versand  durch Arthur Schnitzler am 20. 5. 1906 in Edmund-Weiß-Gasse 7
\newline{}Übermittlung  am 21. 5. 1906 in Wien 18/1 110
\newline{}In Transit  am 21. 5. 1906 in Wien 1/1
\newline{}Zustellung  am 21. 5. 1906 in Wien 1/1
\newline{}Erhalt  durch Gustav Schwarzkopf am 21. 5. 1906 in Tiefer Graben 23}\toendnotes[C]{\smallbreak}
\Standort{CUL, Schnitzler, B 96.}
\physDesc{Postkarte, 181 Zeichen
\newline{}Handschrift: schwarze Tinte, deutsche Kurrent
\newline{}Versand: 1) Rohrpost  2) Stempel: »\nobreak{}\oindex{XVIII., Währing@\textbf{XVIII., Währing}|pwk}Wien {[}18/1{]}
                                       110, 21. \textcolor{gray}{V}. 06, 7\nobreak{}«.  3) Stempel: »\nobreak{}\oindex{I., Innere Stadt@\textbf{I., Innere Stadt}|pwk}Wien 1/1, 21 V 06, 7 10V\nobreak{}«.  4) Stempel: »\nobreak{}\oindex{I., Innere Stadt@\textbf{I., Innere Stadt}|pwk}Wien 1/1, 21 V 06, 7 50V\nobreak{}«. }\toendnotes[C]{\smallbreak}\pstart{}{\pb}\textsc{Hrn
                     Gustav Schwarzkopf}\pend{}\pstart{}Wien I\oindex{I., Innere Stadt@\textbf{I., Innere Stadt}|pw}\pend{}\pstart{}Tiefer Graben 23\oindex{Wien@\textbf{Wien}!I., Innere Stadt@\textbf{I., Innere Stadt}!Tiefer Graben 23@\textbf{Tiefer Graben 23}|pw}.\pend{}{\bigskip}\vspace{1em}
\pstart
           {\pb}\textcolor{gray}{\textbf{Dr. Arthur Schnitzler}}\hfill 20/5. 906\pend
           
\pstart
           \textcolor{gray}{\textbf{Wien,
                        XVIII. Spoettelgasse 7\oindex{XXXX Ortsangabe fehlt|pw}.}}\pend
           \vspace{0.5em}
\pstart
           lieber Guſtav, bitte \label{K_L04037-1v}\edtext{ko{\geminationm}en Sie}{\lemma{\textnormal{\emph{kommen Sie}}}\Cendnote{\textnormal{Schwarzkopf\pwindex{Schwarzkopf, Gustav 7.\,11.\,1853 Wien – 13.\,11.\,1939 ebd.@\textsc{Schwarzkopf, Gustav} (7.\,11.\,1853 Wien – 13.\,11.\,1939 ebd.)|pwk} nahm die
                  Einladung an, siehe Aufführung von Traumulus, 21.5.1906. In: A. S.: \emph{Kulturveranstaltungen}, \url{https://schnitzler-kultur.acdh.oeaw.ac.at/pmb237582.html}.}}}\label{K_L04037-1} morgen \uline{Montag} mit uns zu »Traumulus\pwindex{\textcolor{red}{\textsuperscript{XXXX indx1}}!Traumulus. Tragische Komödie@\strich\emph{Traumulus. Tragische Komödie}|pw}\pwindex{\textcolor{red}{\textsuperscript{XXXX indx1}}!Traumulus. Tragische Komödie@\strich\emph{Traumulus. Tragische Komödie}|pw}\eventindex{Theater an der Wien@\textbf{Theater an der Wien}!Aufführung von Traumulus, 21.5.1906@Aufführung von Traumulus, 21.5.1906|pwv}«.\pend
           
\pstart
           Loge \textsc{parterre} rechts 4.\pend
           
\pstart
           Auf Wiederſehen u herzlichen Gruß{\\[\baselineskip]} Ihr{\\[\baselineskip]}\spacefill\mbox{ArthSch}\pend
           \leftskip=0em{}\selectlanguage{ngerman}\endnumbering\briefempfaengerindex{Schwarzkopf, Gustav@\textsc{Schwarzkopf, Gustav}!zzzSchnitzler, Arthur@\emph{von Arthur Schnitzler}!1906-05-202@{20. 5. 1906}|)be}\mylabel{L04037h}
\begin{anhang}
\end{anhang}\newcommand{\dateiname}{L04037}\newcommand{\titel}{Arthur Schnitzler an Gustav Schwarzkopf, 20. 5. 1906}\newcommand{\editorInnen}{Herausgegeben von Jahnke, SelmaMüller, Martin Anton}\input{../tex-inputs/latex-leseansicht-abspann}
